\chapter{Smooth Manifolds}

\begin{center}\itshape
A sphere is round, but a small piece of it looks flat. A manifold is a space in which every small piece looks like a piece of $\R^n$, glued to its
neighbours smoothly enough that calculus survives the gluing. The inverse function theorem says when a map can be undone locally; it is the one tool from
which all the definitions, all the examples and the first two great theorems of the subject follow.
\end{center}

On 10 June 1854 a shy twenty-seven-year-old gave the trial lecture that would allow him to teach at the University of G\"ottingen. The custom was that
the candidate proposed three topics and the faculty chose the first. Bernhard Riemann's third topic was ``the hypotheses which lie at the foundations of
geometry''. Gauss, who was seventy-seven and had thought about the subject for half a century, chose the third. According to Dedekind, who knew both men,
Gauss came out of the lecture deeply impressed, and on the way home spoke of it to his colleague Wilhelm Weber with the highest praise and an excitement
that was rare for him. The text was published only in 1868, two years after Riemann's death.

The lecture contains almost no formulas. Its key word is \emph{Mannigfaltigkeit}, a ``manifoldness'': a collection of objects each of which is fixed by
$n$ numbers that can vary continuously, like the positions of a point in space, the colours, or the states of a mechanical system. Riemann's proposal was
that the geometry of such a thing is not given in advance by an ambient space; it is given by a rule for measuring lengths of small displacements, which
may change from place to place. That is the curvature of Chapter~4, and sixty years later it was the framework of general relativity.

What a manifold \emph{is} took another eighty years to say precisely. Poincar\'e, in his \emph{Analysis Situs} of 1895, described manifolds of $\R^N$ by
equations and by parametrizations, which are two of the four equivalent descriptions of Section~\ref{sec:submanifolds}. Hermann Weyl, in his 1913 book
on Riemann surfaces, gave the first definition that does not mention any ambient space: a space covered by overlapping neighbourhoods, each identified
with a disc. Hassler Whitney, in 1936, made it the modern definition for every dimension, and proved in the same paper that nothing was lost: every
abstract manifold sits inside some $\R^N$. We prove the easy half of Whitney's theorem at the end of this chapter.

The engine of everything is the inverse function theorem, which we prove from Banach's fixed point theorem. A robot arm is a good example to keep in
mind throughout: its possible positions form a manifold, the points where it loses a degree of freedom are the critical points of a smooth map, and
Sard's theorem, which we prove in the case of equal dimensions, says that such bad positions are rare.

\section{Calculus in \texorpdfstring{$\R^n$}{R\^{}n}, quickly}

We use the Euclidean norm $|x|$ on $\R^n$ and the operator norm $\|A\|=\sup_{|v|\le1}|Av|$ on linear maps. An open set is always non-empty unless we say
otherwise.

\begin{definition}
Let $U\subset\R^n$ be open and $f\colon U\to\R^m$. The map $f$ is \emph{differentiable at $a\in U$} if there is a linear map $Df(a)\colon\R^n\to\R^m$
with
\[
f(a+h)=f(a)+Df(a)h+o(|h|)\qquad(h\to0).
\]
It is \emph{$C^1$} if it is differentiable at every point and $a\mapsto Df(a)$ is continuous, $C^k$ if moreover $Df$ is $C^{k-1}$, and \emph{smooth}
($C^\infty$) if it is $C^k$ for every $k$.
\end{definition}

The matrix of $Df(a)$ is the Jacobian matrix of partial derivatives $\partial_jf_i(a)$, since $h=te_j$ in the definition gives the partial derivative.

\begin{proposition}[chain rule]\label{prop:chain}
If $f$ is differentiable at $a$ and $g$ is differentiable at $f(a)$, then $g\circ f$ is differentiable at $a$ and $D(g\circ f)(a)=Dg(f(a))\,Df(a)$. If
$f$ and $g$ are $C^k$, so is $g\circ f$.
\end{proposition}
\begin{proof}
Write $A=Df(a)$, $B=Dg(b)$ with $b=f(a)$, and $f(a+h)=b+k(h)$ with $k(h)=Ah+\varepsilon_1(h)|h|$, $\varepsilon_1(h)\to0$. Then $|k(h)|\le(\|A\|+1)|h|$ for small
$h$, and
\[
g(b+k)-g(b)-BAh=B\varepsilon_1(h)|h|+\varepsilon_2(k)|k|,
\]
where $\varepsilon_2(k)\to0$ as $k\to0$. Both terms are $o(|h|)$. For the regularity, $D(g\circ f)=(Dg\circ f)\cdot Df$ is a product of continuous
matrix-valued maps, and by induction on $k$, if $f,g$ are $C^k$ then $Dg\circ f$ and $Df$ are $C^{k-1}$, so $D(g\circ f)$ is $C^{k-1}$ (products of
$C^{k-1}$ maps are $C^{k-1}$, by the case of the bilinear product map).
\end{proof}

\begin{proposition}\label{prop:partials}
If all partial derivatives $\partial_jf_i$ exist and are continuous on $U$, then $f$ is $C^1$.
\end{proposition}
\begin{proof}
It suffices to treat $m=1$. Fix $a$, and for small $h$ go from $a$ to $a+h$ along the coordinate directions: with $a^{(0)}=a$ and
$a^{(j)}=a^{(j-1)}+h_je_j$,
\[
f(a+h)-f(a)=\sum_{j=1}^n\bigl(f(a^{(j)})-f(a^{(j-1)})\bigr)=\sum_{j=1}^nh_j\,\partial_jf(c^{(j)})
\]
by the one-variable mean value theorem, with $c^{(j)}$ on the segment $[a^{(j-1)},a^{(j)}]$. Since $c^{(j)}\to a$ as $h\to0$ and $\partial_jf$ is continuous,
$\partial_jf(c^{(j)})=\partial_jf(a)+o(1)$, and the sum is $\sum_jh_j\partial_jf(a)+o(|h|)$.
\end{proof}

The mean value theorem fails for vector-valued maps: $t\mapsto e^{it}$ has $f(2\pi)-f(0)=0$ although $f'$ never vanishes. What survives is an inequality,
and it is all we need.

\begin{lemma}[mean value inequality]\label{lem:mvi}
Let $f\colon U\to\R^m$ be $C^1$ and let the segment $[x,y]$ lie in $U$. If $\|Df\|\le L$ on $[x,y]$, then $|f(y)-f(x)|\le L|y-x|$.
\end{lemma}
\begin{proof}
Let $u$ be a unit vector with $\langle u,f(y)-f(x)\rangle=|f(y)-f(x)|$ and $g(t)=\langle u,f(x+t(y-x))\rangle$ for $t\in[0,1]$. By the chain rule
$g'(t)=\langle u,Df(x+t(y-x))(y-x)\rangle$, so $|g'(t)|\le L|y-x|$, and the one-variable mean value theorem gives
$|f(y)-f(x)|=g(1)-g(0)\le L|y-x|$.
\end{proof}

\begin{lemma}\label{lem:glnopen}
The set $\mathrm{GL}_n(\R)$ of invertible matrices is open in $M_n(\R)$, and inversion $A\mapsto A^{-1}$ is smooth on it. If $\|A-I\|\le\frac12$, then
$A$ is invertible and $\|A^{-1}\|\le2$.
\end{lemma}
\begin{proof}
$\mathrm{GL}_n(\R)=\det^{-1}(\R\setminus\{0\})$ and $\det$ is a polynomial in the entries, so the set is open. By Cramer's rule each entry of $A^{-1}$ is a
polynomial in the entries of $A$ divided by $\det A$, a smooth function on $\mathrm{GL}_n(\R)$. If $\|A-I\|\le\frac12$, then
$|Av|\ge|v|-|(A-I)v|\ge\frac12|v|$, so $A$ is injective, hence invertible, and $|A^{-1}w|\le2|w|$.
\end{proof}

\subsection*{The inverse function theorem}

A linear map $\R^n\to\R^n$ can be undone if and only if it is invertible. The inverse function theorem says that a smooth map can be undone \emph{near a
point} as soon as its linear approximation at that point can.

\begin{definition}
Let $U,V\subset\R^n$ be open. A \emph{$C^k$ diffeomorphism} $f\colon U\to V$ is a bijection such that $f$ and $f^{-1}$ are $C^k$. A \emph{diffeomorphism}
is a $C^\infty$ one.
\end{definition}

\begin{theorem}[inverse function theorem]\label{thm:inverse}
Let $U\subset\R^n$ be open, $k\in\{1,2,\dots,\infty\}$, $f\colon U\to\R^n$ a $C^k$ map and $a\in U$ with $Df(a)$ invertible. There are open sets $V\ni a$
and $W\ni f(a)$ such that $f$ restricts to a $C^k$ diffeomorphism $V\to W$. Its inverse $g$ satisfies $Dg(y)=Df(g(y))^{-1}$.
\end{theorem}
\begin{proof}
\emph{Reduction.} Replacing $f$ by $x\mapsto Df(a)^{-1}\bigl(f(a+x)-f(a)\bigr)$, which changes nothing in the statement, we may assume $a=0$,
$f(0)=0$ and $Df(0)=I$.

\emph{Injectivity on a ball.} By continuity of $Df$, choose $r>0$ with $\bar B(0,2r)\subset U$ and $\|Df(x)-I\|\le\frac12$ for $|x|\le2r$. Let
$\varphi(x)=x-f(x)$. Then $\|D\varphi\|\le\frac12$ on the convex set $\bar B(0,2r)$, so by Lemma~\ref{lem:mvi}, $|\varphi(x)-\varphi(x')|\le\frac12|x-x'|$ there,
and
\[
|f(x)-f(x')|\ge|x-x'|-|\varphi(x)-\varphi(x')|\ge\tfrac12|x-x'|.
\]
In particular $f$ is injective on $\bar B(0,2r)$.

\emph{Surjectivity onto a ball.} Let $|y|<r$ and $T_y(x)=y+\varphi(x)$, so that $T_y(x)=x$ if and only if $f(x)=y$. For $|x|\le2r$,
$|T_y(x)|\le|y|+|\varphi(x)-\varphi(0)|<r+\frac12\cdot2r=2r$. So $T_y$ maps the complete metric space $\bar B(0,2r)$ into itself, and is a
$\frac12$-contraction. By Banach's fixed point theorem (Volume~6, Theorem~2.11) it has a unique fixed point, which lies in the open ball $B(0,2r)$.

Let $W=B(0,r)$ and $V=B(0,2r)\cap f^{-1}(W)$. Both are open, $0\in V$, and $f\colon V\to W$ is a bijection. Let $g\colon W\to V$ be its inverse. The
inequality above, with $x=g(y)$ and $x'=g(y')$, says $|g(y)-g(y')|\le2|y-y'|$, so $g$ is continuous.

\emph{Differentiability.} Fix $y\in W$, let $x=g(y)$ and $A=Df(x)$, invertible with $\|A^{-1}\|\le2$ by Lemma~\ref{lem:glnopen}. For $y+k\in W$ put
$h=g(y+k)-x$, so $|h|\le2|k|$ and $h\to0$ as $k\to0$. Since $f(x+h)=y+k$,
\[
k=Ah+\varepsilon(h)|h|\quad\text{with }\varepsilon(h)\to0,
\]
and therefore
\[
g(y+k)-g(y)-A^{-1}k=h-A^{-1}(Ah+\varepsilon(h)|h|)=-A^{-1}\varepsilon(h)|h|,
\]
whose norm is at most $2|\varepsilon(h)|\cdot2|k|=o(|k|)$. So $g$ is differentiable at $y$ with $Dg(y)=Df(g(y))^{-1}$.

\emph{Regularity.} The formula says $Dg=\mathrm{inv}\circ Df\circ g$, with $\mathrm{inv}$ the smooth inversion map of Lemma~\ref{lem:glnopen}. Since $g$
is continuous and $Df$ is continuous, $Dg$ is continuous and $g$ is $C^1$. If $g$ is $C^j$ for some $j<k$, then $Df\circ g$ is $C^{j}$ (as $Df$ is
$C^{k-1}$ and $j\le k-1$), so $Dg$ is $C^j$ and $g$ is $C^{j+1}$. By induction $g$ is $C^k$.
\end{proof}

\begin{think}
The function $f(x)=x+2x^2\sin\frac1x$, with $f(0)=0$, is differentiable everywhere and $f'(0)=1$. Yet it is injective on no neighbourhood of $0$. Find
points where $f'<0$ arbitrarily close to $0$. Which line of the proof breaks?
\end{think}

The theorem has a twin, which solves equations for some of the variables in terms of the others.

\begin{theorem}[implicit function theorem]\label{thm:implicit}
Let $U\subset\R^n\times\R^m$ be open, $F\colon U\to\R^m$ a $C^k$ map, and $(a,b)\in U$ with $F(a,b)=0$ and with $D_yF(a,b)$, the derivative with respect
to the last $m$ variables, invertible. There are open sets $A\ni a$ in $\R^n$ and $B\ni b$ in $\R^m$ with $A\times B\subset U$, and a $C^k$ map
$h\colon A\to B$, such that for $(x,y)\in A\times B$:
\[
F(x,y)=0\iff y=h(x).
\]
\end{theorem}
\begin{proof}
Let $\Phi(x,y)=(x,F(x,y))$. Its Jacobian at $(a,b)$ is the block matrix $\bigl(\begin{smallmatrix}I&0\\ D_xF&D_yF\end{smallmatrix}\bigr)$, invertible.
By Theorem~\ref{thm:inverse} there are open $V\ni(a,b)$ and $W\ni(a,0)$ with $\Phi\colon V\to W$ a $C^k$ diffeomorphism. Its inverse fixes the first
coordinate, so it has the form $\Psi(x,z)=(x,\psi(x,z))$ with $\psi$ of class $C^k$. Choose open balls $A_0\ni a$ and $B\ni b$ with $A_0\times B\subset V$,
and let
\[
A=\{x\in A_0:\ (x,0)\in W\text{ and }\psi(x,0)\in B\},
\]
an open set containing $a$, since $\psi(a,0)=b$. Put $h(x)=\psi(x,0)$ for $x\in A$; it is $C^k$ with values in $B$, and $F(x,h(x))=0$ because
$\Phi(\Psi(x,0))=(x,0)$. Conversely, if $(x,y)\in A\times B$ and $F(x,y)=0$, then $(x,y)$ and $(x,h(x))$ both lie in $V$ and have the same image $(x,0)$
under $\Phi$, which is injective on $V$, so $y=h(x)$.
\end{proof}

\section{Submanifolds of \texorpdfstring{$\R^N$}{R\^{}N}}\label{sec:submanifolds}

\begin{definition}
A subset $M\subset\R^N$ is a \emph{$k$-dimensional submanifold} if every $p\in M$ has an open neighbourhood $U$ in $\R^N$ with a diffeomorphism $\varphi$
from $U$ onto an open set $V\subset\R^N$ that \emph{straightens} $M$:
\[
\varphi(M\cap U)=V\cap(\R^k\times\{0\}).
\]
\end{definition}

So near each of its points, $M$ is a piece of a $k$-plane seen through a smooth change of coordinates. This is the cleanest definition and the least
practical one: nobody describes the sphere by giving a straightening map. In practice a submanifold comes as the solution set of equations, as a graph,
or as the image of a parametrization, and the following theorem says that all four are the same thing.

\begin{theorem}[four descriptions]\label{thm:fourdesc}
Let $M\subset\R^N$, $p\in M$ and $0\le k\le N$. The following are equivalent.
\begin{enumerate}[label=(\alph*)]
\item \emph{Straightening:} there is a diffeomorphism $\varphi\colon U\to V$ between open sets of $\R^N$, with $p\in U$ and
$\varphi(M\cap U)=V\cap(\R^k\times\{0\})$.
\item \emph{Equations:} there are an open $U\ni p$ and a smooth $F\colon U\to\R^{N-k}$ with $M\cap U=F^{-1}(0)$ and $DF(x)$ surjective for every
$x\in M\cap U$.
\item \emph{Graph:} after a permutation of the coordinates of $\R^N=\R^k\times\R^{N-k}$, there are open sets $\Omega\subset\R^k$ and $B\subset\R^{N-k}$ with
$p\in\Omega\times B$, and a smooth $g\colon\Omega\to B$, such that $M\cap(\Omega\times B)$ is the graph of $g$.
\item \emph{Parametrization:} there are an open $\Omega\subset\R^k$, an open $U\ni p$ and a smooth $\psi\colon\Omega\to\R^N$ such that $\psi$ is injective,
$D\psi(x)$ is injective for every $x$, $\psi(\Omega)=M\cap U$, and $\psi$ is a homeomorphism from $\Omega$ onto $M\cap U$.
\end{enumerate}
So $M$ is a $k$-dimensional submanifold if and only if one of them holds at every point of $M$.
\end{theorem}
\begin{proof}
(a)$\Rightarrow$(b). Let $F$ be the last $N-k$ coordinates of $\varphi$. Then $F^{-1}(0)=\varphi^{-1}(V\cap(\R^k\times\{0\}))=M\cap U$, and the rows of
$DF(x)$ are the last $N-k$ rows of the invertible matrix $D\varphi(x)$, which are independent; so $DF(x)$ has rank $N-k$.

(b)$\Rightarrow$(c). The $(N-k)\times N$ matrix $DF(p)$ has rank $N-k$, so it has $N-k$ independent columns. Permute the coordinates so that they are the
last ones, write points as $(x,y)\in\R^k\times\R^{N-k}$ and $p=(a,b)$; then $D_yF(a,b)$ is invertible. Theorem~\ref{thm:implicit} gives $A\ni a$, $B\ni b$
and $h$; shrinking $A$ and $B$ to balls we may assume $A\times B\subset U$. Then $M\cap(A\times B)=F^{-1}(0)\cap(A\times B)$ is the graph of $h$. Take
$\Omega=A$, $g=h$.

(c)$\Rightarrow$(d). Let $\psi(x)=(x,g(x))$ on $\Omega$, followed by the inverse permutation of coordinates. It is smooth and injective, its derivative
$\bigl(\begin{smallmatrix}I\\ Dg\end{smallmatrix}\bigr)$ is injective, its image is $M\cap(\Omega\times B)$ with $\Omega\times B$ open, and its inverse on
the image is the projection $(x,y)\mapsto x$, which is continuous.

(d)$\Rightarrow$(a). Let $\psi(x_0)=p$. The $N\times k$ matrix $D\psi(x_0)$ has rank $k$, so it has $k$ independent rows; permute the coordinates so that
they are the first $k$. Define $\Theta\colon\Omega\times\R^{N-k}\to\R^N$ by $\Theta(x,z)=\psi(x)+(0,z)$. Its Jacobian at $(x_0,0)$ is
$\bigl(\begin{smallmatrix}D_1\psi&0\\ D_2\psi&I\end{smallmatrix}\bigr)$ with $D_1\psi(x_0)$ invertible, so by Theorem~\ref{thm:inverse} $\Theta$ restricts to a
diffeomorphism from an open $V'\ni(x_0,0)$ onto an open $U'\ni p$. Note that $\Theta(x,0)=\psi(x)$.

Here the homeomorphism condition enters. The set $\Omega_0=\{x:(x,0)\in V'\}$ is open in $\Omega$, so $\psi(\Omega_0)$ is open in $M\cap U$, that is,
$\psi(\Omega_0)=M\cap O$ for some open $O\subset U$. Let $U''=U'\cap O$ and $V''=V'\cap\Theta^{-1}(U'')$. We claim
\[
\Theta\bigl(V''\cap(\R^k\times\{0\})\bigr)=M\cap U''.
\]
If $(x,0)\in V''$, then $\Theta(x,0)=\psi(x)$ lies in $M$ and in $U''$. Conversely, if $q\in M\cap U''$, then $q\in M\cap O=\psi(\Omega_0)$, so $q=\psi(x)=\Theta(x,0)$
with $(x,0)\in V'$, and $\Theta(x,0)\in U''$ gives $(x,0)\in V''$. So $\varphi=(\Theta|_{V''})^{-1}\colon U''\to V''$ straightens $M$ near $p$.
\end{proof}

\begin{think}
The curve $\psi(t)=(\sin2t,\sin t)$ on $\Omega=\left]0,2\pi\right[$ is injective and its derivative never vanishes, but its image, a figure eight, is not a
submanifold at the origin (Exercise~\ref{ex:nosubmanifold}). Which hypothesis of (d) fails, and where in the proof was it used?
\end{think}

The number $k$ is determined by $M$: we will see in Section~\ref{sec:tangent} that it is the dimension of a vector space attached to each point. A
$0$-dimensional submanifold is a set of isolated points, an $N$-dimensional one is an open set.

\begin{example}[spheres]\label{ex:spheres}
$S^n=\{x\in\R^{n+1}:|x|=1\}$ is the zero set of $F(x)=|x|^2-1$, and $DF(x)h=2\langle x,h\rangle$ is a non-zero linear form when $x\ne0$. By (b), $S^n$ is an
$n$-dimensional submanifold of $\R^{n+1}$.
\end{example}

\begin{example}[the orthogonal group]\label{ex:on}
Let $\mathrm{Sym}_n$ be the space of symmetric $n\times n$ matrices, of dimension $\frac{n(n+1)}2$, and $F\colon M_n(\R)\to\mathrm{Sym}_n$,
$F(X)=X^\top X-I$, so that $\mathrm O(n)=F^{-1}(0)$. Expanding $F(X+H)$ gives
\[
DF(X)H=H^\top X+X^\top H.
\]
If $X\in\mathrm O(n)$ and $S\in\mathrm{Sym}_n$, then $H=\frac12XS$ satisfies $H^\top X=\frac12SX^\top X=\frac12S$ and $X^\top H=\frac12S$, so
$DF(X)H=S$: the derivative is surjective onto $\mathrm{Sym}_n$. By (b), $\mathrm O(n)$ is a submanifold of $M_n(\R)=\R^{n^2}$ of dimension
\[
n^2-\tfrac{n(n+1)}2=\tfrac{n(n-1)}2.
\]
It is compact, being closed and bounded (each column is a unit vector). The determinant is $\pm1$ on $\mathrm O(n)$, so $\mathrm{SO}(n)=\mathrm O(n)\cap\{\det>0\}$
is open in $\mathrm O(n)$, hence also a submanifold of dimension $\frac{n(n-1)}2$. For $n=2$ this is a circle; for $n=3$ it is three-dimensional, which
matches the count ``an axis (two parameters) and an angle (one)''.
\end{example}

The rotations of space form a smooth three-dimensional manifold; Chapter~7 will show it is $\R P^3$, and that the quaternions of Volume~2 are its double
cover.

\begin{example}[Cassini ovals]\label{ex:cassini}
Let $f(x,y)=(x^2+y^2)^2-2(x^2-y^2)$. Writing $z=x+iy$ one checks $f=|z-1|^2|z+1|^2-1$, so the level set $f=c$ is the set of points whose distances to
$\pm1$ have product $\sqrt{1+c}$: a Cassini oval. The gradient is
\[
\nabla f=\bigl(4x(x^2+y^2-1),\ 4y(x^2+y^2+1)\bigr),
\]
which vanishes only at $(0,0)$ and $(\pm1,0)$, with values $0$ and $-1$. So every level $c>-1$ other than $0$ is a curve, by (b); for
$-1<c<0$ it has two components, for $c>0$ one (Figure~\ref{fig:cassini}). The level $0$ is Jacob Bernoulli's lemniscate, a figure eight, which is not a
submanifold at the origin; the level $-1$ is the two points $(\pm1,0)$, which \emph{is} a $0$-dimensional submanifold: the condition in (b) is
sufficient, not necessary.
\end{example}

\begin{figure}[ht]
\centering
\fig{ch01-cassini}
\caption{Level sets of $f(x,y)=(x^2+y^2)^2-2(x^2-y^2)$. Below the critical value $0$ (charcoal) each level is two ovals around the minima at $(\pm1,0)$; above
it (ochre) one oval. Exactly at the critical value the level set is the lemniscate (coral), which has a crossing at the saddle point $(0,0)$ and is not
a manifold there. The marked points are the three critical points, where $f$ takes the critical values $-1$, $0$ and $-1$.}
\label{fig:cassini}
\end{figure}

\section{Abstract manifolds}\label{sec:abstract}

The torus as a surface of revolution in $\R^3$ is a submanifold, but so is the flat torus $S^1\times S^1\subset\R^4$, and they are ``the same''. The
space of lines through the origin of $\R^3$ is clearly a two-dimensional object, but it is not given to us inside any $\R^N$. Weyl's insight was to
build a manifold from its charts.

\begin{definition}
A \emph{topological manifold of dimension $n$} is a topological space $M$ that is Hausdorff, has a countable basis of open sets, and in which every point
has an open neighbourhood $U$ homeomorphic to an open subset of $\R^n$. Such a homeomorphism $\varphi\colon U\to\varphi(U)\subset\R^n$ is a
\emph{chart}. Two charts $(U,\varphi)$, $(U',\varphi')$ are \emph{compatible} if the \emph{transition map}
\[
\varphi'\circ\varphi^{-1}\colon\varphi(U\cap U')\to\varphi'(U\cap U')
\]
is smooth, and so is its inverse $\varphi\circ\varphi'^{-1}$. An \emph{atlas} is a family of pairwise compatible charts whose domains cover $M$.
\end{definition}

\begin{proposition}\label{prop:maxatlas}
Every atlas $\mathcal A$ is contained in a unique maximal atlas, the set $\bar{\mathcal A}$ of all charts compatible with every chart of $\mathcal A$.
\end{proposition}
\begin{proof}
We must show that two charts $(U,\varphi)$, $(U',\varphi')$ of $\bar{\mathcal A}$ are compatible with each other. Smoothness is local, so fix
$x=\varphi(q)$ with $q\in U\cap U'$, and a chart $(U_0,\varphi_0)\in\mathcal A$ at $q$. Near $x$,
\[
\varphi'\circ\varphi^{-1}=(\varphi'\circ\varphi_0^{-1})\circ(\varphi_0\circ\varphi^{-1}),
\]
a composite of smooth maps. Any atlas containing $\mathcal A$ consists of charts compatible with $\mathcal A$, so it lies in $\bar{\mathcal A}$, which
is therefore maximal and unique.
\end{proof}

\begin{definition}
A \emph{smooth manifold} is a topological manifold with a maximal atlas, its \emph{smooth structure}; to give one it suffices to give any atlas. A map
$f\colon M\to N$ between smooth manifolds is \emph{smooth} if it is continuous and $\psi\circ f\circ\varphi^{-1}$ is smooth for all charts $\varphi$ of $M$
and $\psi$ of $N$, where defined. A \emph{diffeomorphism} is a smooth bijection with a smooth inverse.
\end{definition}

Because the transition maps are smooth, it is enough to check smoothness of $f$ for one chart around each point of $M$ and one around its image. A chart
$\varphi\colon U\to\varphi(U)$ is itself a diffeomorphism, when $U$ and $\varphi(U)$ carry the obvious structures.

\begin{proposition}\label{prop:subisman}
A $k$-dimensional submanifold $M\subset\R^N$, with the subspace topology, is a smooth manifold: for each straightening map $\varphi\colon U\to V$, the
map $\pi\circ\varphi|_{M\cap U}$, where $\pi(x_1,\dots,x_N)=(x_1,\dots,x_k)$, is a chart, and these charts form an atlas. For this structure, the
restriction to $M$ of a smooth map defined on an open set of $\R^N$ is smooth.
\end{proposition}
\begin{proof}
Hausdorff and a countable basis are inherited from $\R^N$. The map $\pi\circ\varphi$ is a homeomorphism from $M\cap U$ onto the open set
$\{x\in\R^k:(x,0)\in V\}$, with inverse $x\mapsto\varphi^{-1}(x,0)$. For two straightenings $\varphi$, $\varphi'$, the transition map is
$x\mapsto\pi\varphi'\varphi^{-1}(x,0)$, a restriction of the smooth map $\pi\circ\varphi'\circ\varphi^{-1}$ to an open part of $\R^k\times\{0\}$, hence
smooth. Finally, if $f$ is smooth on an open set of $\R^N$, then in the chart it reads $x\mapsto f(\varphi^{-1}(x,0))$, which is smooth.
\end{proof}

\begin{example}[products and tori]
If $M$ and $N$ are smooth manifolds of dimensions $m$ and $n$, the charts $\varphi\times\psi$ form an atlas on $M\times N$, of dimension $m+n$: transition
maps are products of transition maps. The $n$-torus $T^n=S^1\times\dots\times S^1$ is a smooth manifold of dimension $n$, and also a submanifold of
$\R^{2n}$, by (b) applied to the $n$ equations $x_{2j-1}^2+x_{2j}^2=1$.
\end{example}

\begin{example}[projective space]\label{ex:rpn}
Let $\R P^n$ be the set of lines through $0$ in $\R^{n+1}$, the quotient of $\R^{n+1}\setminus\{0\}$ by $x\sim\lambda x$ ($\lambda\ne0$), with the quotient
topology; write $[x]=[x_0:\dots:x_n]$. The quotient map $\pi$ is open, because $\pi^{-1}(\pi(O))=\bigcup_{\lambda\ne0}\lambda O$ is open when $O$ is.
\begin{itemize}
\item \emph{Charts.} On the open set $U_i=\{[x]:x_i\ne0\}$ let $\varphi_i([x])=\bigl(\frac{x_j}{x_i}\bigr)_{j\ne i}\in\R^n$. It is a homeomorphism onto
$\R^n$, with inverse $u\mapsto[u_0:\dots:1:\dots:u_n]$ ($1$ in position $i$). On $\varphi_i(U_i\cap U_j)$ the transition map $\varphi_j\circ\varphi_i^{-1}$
divides every coordinate by $u_j\ne0$, after inserting a $1$ in position $i$: it is smooth.
\item \emph{Countable basis.} $\R P^n$ is the union of $n+1$ open sets homeomorphic to $\R^n$.
\item \emph{Hausdorff.} The map $P([x])=\frac{xx^\top}{|x|^2}$ is well defined and continuous on $\R P^n$, with values in the matrices. It is injective,
because $P([x])$ is the orthogonal projection onto the line $\R x$, which determines the line. Two points with distinct images under a continuous
injective map into a Hausdorff space have disjoint neighbourhoods.
\end{itemize}
Moreover $\R P^n=\pi(S^n)$ is compact.
\end{example}

\begin{remark}
The two topological conditions are not decoration. The \emph{line with two origins}, two copies of $\R$ glued along $\R\setminus\{0\}$, has an atlas of two
charts with identity transition maps, but its two origins cannot be separated: it is not Hausdorff, and a sequence $\frac1k$ has two limits.
Second-countability excludes monsters such as the long line, and it is what makes partitions of unity, below, countable.
\end{remark}

\subsection*{Bump functions and partitions of unity}

Analytic functions are rigid: one that vanishes on an open set vanishes everywhere. Smooth functions are not, and this flexibility is what makes smooth
manifolds pleasant to work with.

\begin{lemma}\label{lem:bump}
The function $h(t)=e^{-1/t}$ for $t>0$, $h(t)=0$ for $t\le0$, is smooth on $\R$. Hence $\beta(t)=\frac{h(t)}{h(t)+h(1-t)}$ is a smooth function with
$\beta=0$ on $\left]-\infty,0\right]$ and $\beta=1$ on $[1,+\infty[$, and $b(x)=\beta(2-|x|)$ is a smooth function on $\R^n$ with $0\le b\le1$, $b=1$ on
$\bar B(0,1)$ and $b=0$ outside $B(0,2)$.
\end{lemma}
\begin{proof}
For $t>0$, induction shows $h^{(k)}(t)=P_k(\frac1t)e^{-1/t}$ for a polynomial $P_k$, which tends to $0$ as $t\to0^+$ since $e^{-s}$ beats every power of $s$.
Suppose $h^{(k)}$ exists on $\R$, is continuous, and vanishes on $\left]-\infty,0\right]$. Then $h^{(k)}(t)/t\to0$ as $t\to0^\pm$, so $h^{(k+1)}(0)=0$ exists and
$h^{(k+1)}$ is continuous at $0$. The denominator of $\beta$ never vanishes, since $t$ and $1-t$ are not both $\le0$. Finally $b$ is smooth away from $0$,
where $|x|$ is smooth, and equal to $1$ on $B(0,1)$.
\end{proof}

\begin{theorem}[partitions of unity]\label{thm:partition}
Let $M$ be a smooth manifold and $(O_i)_{i\in I}$ an open cover. There is a countable family of smooth functions $\rho_j\colon M\to[0,1]$ such that each
support $\operatorname{supp}\rho_j=\overline{\{\rho_j\ne0\}}$ is compact and contained in some $O_i$, each point has a neighbourhood meeting only finitely
many supports, and $\sum_j\rho_j=1$.
\end{theorem}
\begin{proof}
\emph{Precompact basis.} Each point $p$ has a chart $\varphi$ with $\varphi(p)=0$, and for small $r$ the set $\varphi^{-1}(B(0,r))$ has compact closure
$\varphi^{-1}(\bar B(0,r))$ (compact sets are closed in a Hausdorff space). The members of a countable basis that lie in such a set are still a basis,
so $M$ has a countable basis $B_1,B_2,\dots$ of open sets with compact closures.

\emph{Exhaustion.} Let $K_1=\bar B_1$. Given the compact $K_m$, it is covered by finitely many $B_l$; let $j\ge m+1$ be such that $K_m\subset B_1\cup\dots\cup B_j$
and let $K_{m+1}=\bar B_1\cup\dots\cup\bar B_j$. Then the $K_m$ are compact, $K_m\subset\operatorname{int}K_{m+1}$, and $\bigcup_mK_m=M$ since
$B_m\subset K_m$. Put $K_0=K_{-1}=\emptyset$.

\emph{Bumps.} The set $A_m=K_m\setminus\operatorname{int}K_{m-1}$ is compact, and contained in the open set $W_m=\operatorname{int}K_{m+1}\setminus K_{m-2}$.
For each $p\in A_m$ choose $i$ with $p\in O_i$, a chart $\varphi$ with $\varphi(p)=0$, and $r>0$ with $\varphi^{-1}(\bar B(0,2r))\subset O_i\cap W_m$. The
function $\sigma_p=b(\varphi/r)$ on the chart domain, extended by $0$, is smooth on $M$: its support $\varphi^{-1}(\bar B(0,2r))$ is compact, hence closed,
and off it $\sigma_p$ vanishes on an open set. It equals $1$ on $\varphi^{-1}(B(0,r))$, and finitely many of these sets cover $A_m$. Collect the
corresponding $\sigma_p$ for all $m$: a countable family $(\sigma_j)$.

\emph{Normalization.} A point of $\operatorname{int}K_m$ has a neighbourhood that misses $W_{m'}$ for $m'\ge m+2$, hence meets only the finitely many
supports coming from $m'\le m+1$. So $\sigma=\sum_j\sigma_j$ is locally a finite sum of smooth functions, hence smooth, and $\sigma\ge1$ everywhere since
the sets $A_m$ cover $M$. Put $\rho_j=\sigma_j/\sigma$.
\end{proof}

\section{Tangent spaces}\label{sec:tangent}

For a submanifold there is an obvious candidate: the set of velocities of curves drawn on it.

\begin{proposition}\label{prop:tangentsub}
Let $M\subset\R^N$ be a $k$-dimensional submanifold and $p\in M$. Let $T_pM$ be the set of velocities $\gamma'(0)$ of smooth curves
$\gamma\colon\left]-\varepsilon,\varepsilon\right[\to\R^N$ with values in $M$ and $\gamma(0)=p$. For any $F$ as in Theorem~\ref{thm:fourdesc}(b) and any $\psi$
as in (d) with $\psi(x_0)=p$,
\[
T_pM=\operatorname{im}D\psi(x_0)=\ker DF(p).
\]
In particular $T_pM$ is a $k$-dimensional vector space, and $k$ is determined by $M$.
\end{proposition}
\begin{proof}
If $v\in\R^k$, the curve $t\mapsto\psi(x_0+tv)$ lies in $M$ and has velocity $D\psi(x_0)v$, so $\operatorname{im}D\psi(x_0)\subset T_pM$. If $\gamma$ is a curve in $M$,
then $F(\gamma(t))=0$ near $t=0$, and differentiating gives $DF(p)\gamma'(0)=0$, so $T_pM\subset\ker DF(p)$. The first space has dimension $k$ because
$D\psi(x_0)$ is injective, the last has dimension $N-(N-k)=k$ because $DF(p)$ is surjective, so the three coincide.
\end{proof}

For $S^n$ this gives $T_pS^n=p^\perp$, and for $\mathrm O(n)$ at the identity, $T_I\mathrm O(n)=\{H:H^\top+H=0\}$, the skew-symmetric matrices. These are
the first examples of Lie algebras, the subject of Chapter~7.

On an abstract manifold there is nothing for a curve to have a velocity in. The way out is to remember what a tangent vector \emph{does}: it
differentiates functions. If $v=\gamma'(0)$, then $f\mapsto(f\circ\gamma)'(0)$ is linear and obeys the Leibniz rule.

\begin{definition}
Let $M$ be a smooth manifold and $C^\infty(M)$ the algebra of smooth functions $M\to\R$. A \emph{tangent vector at $p$} is a linear map
$v\colon C^\infty(M)\to\R$ with $v(fg)=f(p)\,v(g)+g(p)\,v(f)$. They form a vector space $T_pM$, the \emph{tangent space}.
\end{definition}

\begin{lemma}\label{lem:derivlocal}
Let $v\in T_pM$. Then $v$ kills constants, and if $f=g$ on a neighbourhood of $p$, then $v(f)=v(g)$.
\end{lemma}
\begin{proof}
$v(1)=v(1\cdot1)=2v(1)$, so $v(1)=0$. If $f=g$ on an open $O\ni p$, take a smooth $\chi$ with $\chi(p)=1$ and support in $O$ (a bump in a chart, as in the
proof of Theorem~\ref{thm:partition}). Then $\chi(f-g)=0$, so $0=v(\chi(f-g))=\chi(p)v(f-g)+(f-g)(p)v(\chi)=v(f-g)$.
\end{proof}

\begin{lemma}[Hadamard]\label{lem:hadamard}
Let $f$ be smooth on a ball $B\subset\R^n$ centred at $a$. There are smooth $g_1,\dots,g_n$ on $B$ with $g_i(a)=\partial_if(a)$ and
$f(x)=f(a)+\sum_i(x_i-a_i)g_i(x)$.
\end{lemma}
\begin{proof}
$f(x)-f(a)=\int_0^1\frac{\dd}{\dd t}f(a+t(x-a))\dd t=\sum_i(x_i-a_i)\int_0^1\partial_if(a+t(x-a))\dd t$. Put $g_i(x)=\int_0^1\partial_if(a+t(x-a))\dd t$; it is smooth
by differentiation under the integral sign, every derivative of the integrand being continuous on the compact $[0,1]\times\bar B'$ for slightly smaller
closed balls $\bar B'$.
\end{proof}

\begin{theorem}\label{thm:tangentbasis}
Let $(U,\varphi)$ be a chart at $p$ with coordinates $\varphi=(x^1,\dots,x^n)$. The maps
\[
\frac{\partial}{\partial x^i}\Big|_p\colon f\mapsto\partial_i(f\circ\varphi^{-1})(\varphi(p))
\]
form a basis of $T_pM$. In particular $\dim T_pM=n$.
\end{theorem}
\begin{proof}
They are tangent vectors by the Leibniz rule for partial derivatives. Translate so that $\varphi(p)=0$, choose $\varepsilon$ with
$B(0,\varepsilon)\subset\varphi(U)$, and a smooth $\chi$ on $M$, equal to $1$ near $p$, whose support is a compact subset of $\varphi^{-1}(B(0,\varepsilon))$
(Lemma~\ref{lem:bump}: take $\chi=b(4\varphi/\varepsilon)$). The functions $X^j=\chi x^j$, extended by $0$, are smooth on $M$ and agree with $x^j$ near $p$, so by
Lemma~\ref{lem:derivlocal} $\frac{\partial}{\partial x^i}|_p(X^j)=\delta_{ij}$: the $n$ vectors are independent.

Let $v\in T_pM$ and $f\in C^\infty(M)$. Lemma~\ref{lem:hadamard} on $B(0,\varepsilon)$ gives $f\circ\varphi^{-1}(y)=f(p)+\sum_iy_ig_i(y)$ with
$g_i(0)=\partial_i(f\circ\varphi^{-1})(0)$. The functions $G_i=\chi\cdot(g_i\circ\varphi)$, extended by $0$, are smooth on $M$, and near $p$ we have
$f=f(p)+\sum_iX^iG_i$. By Lemma~\ref{lem:derivlocal} and the Leibniz rule,
\[
v(f)=\sum_i\bigl(v(X^i)G_i(p)+X^i(p)v(G_i)\bigr)=\sum_iv(X^i)\,\frac{\partial f}{\partial x^i}\Big|_p,
\]
since $X^i(p)=0$. So $v=\sum_iv(X^i)\frac{\partial}{\partial x^i}|_p$.
\end{proof}

\begin{definition}
The \emph{differential} of a smooth $F\colon M\to N$ at $p$ is the linear map $dF_p\colon T_pM\to T_{F(p)}N$, $dF_p(v)(g)=v(g\circ F)$.
\end{definition}

The chain rule $d(G\circ F)_p=dG_{F(p)}\circ dF_p$ is now a tautology. In charts $(x^i)$ at $p$ and $(y^j)$ at $F(p)$, the matrix of $dF_p$ in the bases of
Theorem~\ref{thm:tangentbasis} is the Jacobian matrix of $y\circ F\circ x^{-1}$, by the chain rule of Proposition~\ref{prop:chain}; so a diffeomorphism has
invertible differentials, and $dF_p$ is the abstract face of the derivative.

The two definitions of the tangent space agree for submanifolds. If $M\subset\R^N$ and $w=\gamma'(0)$ is a velocity, then $f\mapsto(f\circ\gamma)'(0)$ is a
derivation that depends only on $w$: writing $\gamma=\psi\circ c$ near $0$ with $c=\psi^{-1}\circ\gamma$ (smooth, since $\psi^{-1}=\pi\circ\varphi$ is the
restriction of a smooth map), $(f\circ\gamma)'(0)=D(f\circ\psi)(x_0)\,D\psi(x_0)^{-1}w$. This is linear in $w$, and injective because the coordinate
functions of $\R^N$, restricted to $M$, recover $w$. Both spaces have dimension $k$, so the map is an isomorphism, and we will use whichever picture is
more convenient.

\section{The regular value theorem}

\begin{definition}
Let $F\colon M\to N$ be smooth. A point $p$ is \emph{critical} if $dF_p$ is not surjective, and \emph{regular} otherwise; a point $q\in N$ is a
\emph{critical value} if $F^{-1}(q)$ contains a critical point, and a \emph{regular value} otherwise (including when $F^{-1}(q)$ is empty). The map $F$ is
an \emph{immersion} if every $dF_p$ is injective, a \emph{submersion} if every $dF_p$ is surjective, and an \emph{embedding} if it is an immersion and a
homeomorphism onto its image.
\end{definition}

\begin{definition}
A subset $S$ of an $n$-dimensional manifold $M$ is a \emph{$k$-dimensional embedded submanifold} if every $p\in S$ is in the domain of a chart $(U,\varphi)$
of $M$ with $\varphi(S\cap U)=\varphi(U)\cap(\R^k\times\{0\})$.
\end{definition}

Exactly as in Proposition~\ref{prop:subisman}, the maps $\pi\circ\varphi|_{S\cap U}$ make $S$ a smooth manifold, the inclusion $\iota\colon S\to M$ is smooth, and
$d\iota_p$ is injective (in these charts $\iota$ reads $x\mapsto(x,0)$). For $M=\R^N$ we recover the submanifolds of Section~\ref{sec:submanifolds}.

\begin{theorem}[regular value theorem]\label{thm:regularvalue}
Let $F\colon M\to N$ be smooth, $\dim M=m$, $\dim N=n$, and let $q$ be a regular value with $S=F^{-1}(q)$ non-empty. Then $S$ is an embedded submanifold of
$M$ of dimension $m-n$, and $T_pS=\ker dF_p$ for every $p\in S$.
\end{theorem}
\begin{proof}
Let $p\in S$, and choose charts $\varphi$ at $p$ and $\psi$ at $q$ with $\varphi(p)=0$, $\psi(q)=0$, and $F$ mapping the domain of $\varphi$ into that of $\psi$.
Then $G=\psi\circ F\circ\varphi^{-1}$ is a smooth map from an open set of $\R^m$ to $\R^n$, and $\varphi(S\cap U)=G^{-1}(0)$. The Jacobian $DG(x)$ is surjective
at every point of $G^{-1}(0)$, because these points come from regular points of $F$. By (b)$\Rightarrow$(a) of Theorem~\ref{thm:fourdesc} there is a diffeomorphism
$\theta$ near $0$ in $\R^m$ with $\theta(G^{-1}(0))=\R^{m-n}\times\{0\}$ locally, and $\theta\circ\varphi$ is the required chart of $M$.

Since $F\circ\iota$ is constant on $S$, $dF_p\circ d\iota_p=0$, so $d\iota_p(T_pS)\subset\ker dF_p$. Both have dimension $m-n$ (the kernel because $dF_p$ is
surjective), so they are equal.
\end{proof}

\begin{proposition}\label{prop:embedding}
If $F\colon M\to N$ is an embedding, then $F(M)$ is an embedded submanifold of $N$, and $F$ is a diffeomorphism from $M$ onto $F(M)$. If $M$ is compact,
every injective immersion is an embedding.
\end{proposition}
\begin{proof}
Work near $p\in M$ in charts $\varphi$ at $p$ and $\psi$ at $F(p)$; then $\psi\circ F\circ\varphi^{-1}$ is an injective immersion of an open set of $\R^m$, and
a homeomorphism onto its image, which is open in $\psi(F(M)\cap V)$ because $F$ is a homeomorphism onto $F(M)$. The proof of (d)$\Rightarrow$(a) in
Theorem~\ref{thm:fourdesc} applies verbatim and produces the straightening chart. In that chart and in $\varphi$, $F$ reads $x\mapsto(x,0)$, so $F\colon M\to F(M)$ has a smooth inverse. If $M$ is compact, a continuous injection into the Hausdorff space $N$ maps closed
(hence compact) sets to compact (hence closed) sets, so it is a homeomorphism onto its image.
\end{proof}

\section{Sard's theorem in equal dimensions}

Critical points can be plentiful: a constant map has only critical points. Critical \emph{values} cannot. Arthur Sard proved in 1942 that the critical
values of a smooth map always form a set of measure zero; we prove the case of equal dimensions, which is the one Chapter~4 will use, and which needs
only $C^1$.

\begin{lemma}\label{lem:nullsets}
Let $U\subset\R^n$ be open and $f\colon U\to\R^n$ a $C^1$ map. If $Z\subset U$ has Lebesgue measure zero, so does $f(Z)$.
\end{lemma}
\begin{proof}
$U$ is a countable union of closed cubes, so it suffices to treat $Z\cap Q$ for a closed cube $Q\subset U$. On the compact convex $Q$, $\|Df\|\le L$ for
some $L$, so $f$ is $L$-Lipschitz there by Lemma~\ref{lem:mvi}. Given $\varepsilon>0$, cover $Z\cap Q$ by countably many cubes $C_j$ of sides $s_j$ with
$\sum_js_j^n<\varepsilon$ (Volume~6, Theorem~6.11). The set $f(C_j\cap Q)$ has diameter at most $L\sqrt n\,s_j$, so it lies in a cube of volume
$(2L\sqrt n)^ns_j^n$, and $f(Z\cap Q)$ is covered by cubes of total volume at most $(2L\sqrt n)^n\varepsilon$.
\end{proof}

\begin{theorem}[Sard, equal dimensions]\label{thm:sard}
Let $U\subset\R^n$ be open and $f\colon U\to\R^n$ a $C^1$ map, and let $C=\{x\in U:\det Df(x)=0\}$ be its set of critical points. Then $f(C)$ has
measure zero.
\end{theorem}
\begin{proof}
Again it suffices to show that $f(C\cap Q)$ is null for a closed cube $Q\subset U$, of side $\ell$. Let $M=\max_Q\|Df\|$. Fix $\varepsilon>0$. By Heine's
theorem (Volume~6, Theorem~4.9) $Df$ is uniformly continuous on $Q$: there is $\delta>0$ with $\|Df(x)-Df(y)\|\le\varepsilon$ when $|x-y|\le\delta$. Cut
$Q$ into $k^n$ subcubes of side $\frac\ell k$, with $k$ so large that their diameter $d=\frac{\sqrt n\,\ell}k$ is at most $\delta$.

Let $Q'$ be a subcube containing a critical point $c$. For $x\in Q'$, the map $x\mapsto f(x)-Df(c)x$ has derivative of norm at most $\varepsilon$ on
$Q'$, so by Lemma~\ref{lem:mvi}
\[
\bigl|f(x)-f(c)-Df(c)(x-c)\bigr|\le\varepsilon|x-c|\le\varepsilon d.
\]
Since $Df(c)$ is singular, its image lies in a hyperplane $H$. So $f(Q')$ lies within distance $\varepsilon d$ of the set $f(c)+\{w\in H:|w|\le Md\}$. In an
orthonormal basis whose last vector is normal to $H$, this neighbourhood is contained in a box with $n-1$ sides of length $2(M+\varepsilon)d$ and one side
of length $2\varepsilon d$. Lebesgue measure is invariant under rotations (Volume~6, Theorem~8.7), so
\[
\lambda^*\bigl(f(Q')\bigr)\le2^n(M+\varepsilon)^{n-1}\varepsilon\,d^n.
\]
There are at most $k^n$ subcubes, and $k^nd^n=n^{n/2}\ell^n$, so
\[
\lambda^*\bigl(f(C\cap Q)\bigr)\le2^n(M+\varepsilon)^{n-1}n^{n/2}\ell^n\,\varepsilon.
\]
This holds for every $\varepsilon>0$, so $f(C\cap Q)$ is null.
\end{proof}

\begin{beauty}
The proof is one picture: near a critical point the map squashes a small cube of size $d$ into a pancake of thickness $\varepsilon d$, so the image has
volume $\varepsilon d^n$ instead of $d^n$. Summing over the cubes, the critical values fill a fraction $\varepsilon$ of the space, for every $\varepsilon$.
The critical points themselves can fill a whole open set (for a constant map they are everything); what is small is always the set of values.
\end{beauty}

On a manifold, a set $Z\subset N$ has \emph{measure zero} if $\psi(Z\cap V)$ is Lebesgue-null for every chart $(V,\psi)$. By Lemma~\ref{lem:nullsets}
applied to the transition maps, it suffices to check this on the charts of one atlas.

\begin{corollary}\label{cor:sardmfd}
Let $F\colon M\to N$ be a smooth map between manifolds of the same dimension. Its set of critical values has measure zero; in particular its regular
values are dense in $N$.
\end{corollary}
\begin{proof}
Since $M$ has a countable basis, countably many charts $(U_i,\varphi_i)$ cover it: the members of the basis that lie in some chart domain
cover $M$, and for each of them we pick one such chart. Fix a chart $(V,\psi)$ of $N$. In the coordinates $G_i=\psi\circ F\circ\varphi_i^{-1}$ on $\varphi_i(U_i\cap F^{-1}(V))$, a point is
critical for $F$ exactly when it is critical for $G_i$, because charts have invertible differentials. So $\psi(\text{critical values}\cap V)$ is the union over
$i$ of the critical values of the $G_i$, a countable union of null sets by Theorem~\ref{thm:sard}. A null set contains no non-empty open set of
$\R^n$, so its complement is dense.
\end{proof}

\begin{remark}
The full theorem of Sard (1942) says that for a $C^r$ map $\R^m\to\R^n$ with $r\ge\max(m-n+1,1)$ the critical values are null; A.~P.~Morse had proved the
case $n=1$ in 1939. The smoothness hypothesis cannot be weakened to $C^1$ when $m>n$: Whitney constructed in 1935 a $C^1$ function on $\R^2$ which is not
constant on a connected arc of critical points, so that its critical values fill an interval. When $m<n$ every point is critical and the theorem says
that the image is null, which Exercise~\ref{ex:minisard} proves directly.
\end{remark}

\section{Every compact manifold lives in some \texorpdfstring{$\R^N$}{R\^{}N}}

\begin{theorem}[Whitney, easy version]\label{thm:whitney}
Every compact smooth manifold $M$ admits an embedding into some $\R^N$. Consequently $M$ is diffeomorphic to a submanifold of $\R^N$.
\end{theorem}
\begin{proof}
Let $n=\dim M$. For each $p\in M$ choose a chart $(U,\varphi)$ with $\varphi(p)=0$, $\bar B(0,2r)\subset\varphi(U)$, and the bump $\rho=b(\varphi/r)$ of
Lemma~\ref{lem:bump}, extended by $0$: it is smooth on $M$, its support is a compact subset of $U$, and it equals $1$ on the open set
$W=\varphi^{-1}(B(0,r))\ni p$. By compactness finitely many of the $W$ cover $M$; call the corresponding data $(U_i,\varphi_i,\rho_i,W_i)$, $1\le i\le s$.
The functions $\rho_i\varphi_i$, extended by $0$ outside $U_i$, are smooth on $M$ (smooth on $U_i$, and zero on the open set $M\setminus\operatorname{supp}\rho_i$;
these two sets cover $M$). Define
\[
F=(\rho_1\varphi_1,\dots,\rho_s\varphi_s,\rho_1,\dots,\rho_s)\colon M\to\R^{sn+s}.
\]
\emph{Injective.} If $F(p)=F(q)$, pick $i$ with $p\in W_i$. Then $\rho_i(q)=\rho_i(p)=1$, so $q\in U_i$, and
$\varphi_i(q)=\rho_i\varphi_i(q)=\rho_i\varphi_i(p)=\varphi_i(p)$; since $\varphi_i$ is injective, $p=q$.

\emph{Immersion.} On the open set $W_i$ the component $\rho_i\varphi_i$ equals the chart $\varphi_i$, whose differential is an isomorphism. So for
$p\in W_i$, $dF_p$ is injective already on these $n$ coordinates.

By Proposition~\ref{prop:embedding}, $F$ is an embedding and a diffeomorphism onto the submanifold $F(M)$.
\end{proof}

\begin{remark}
Whitney's 1936 paper proves much more: every $n$-dimensional manifold, compact or not, embeds in $\R^{2n+1}$. The idea is to start from the embedding above
and project it to hyperplanes one direction at a time; Sard's theorem, applied to suitable maps into the sphere of directions, shows that almost every
direction can be projected away while the map stays an injective immersion, as long as $N>2n+1$. In 1944 Whitney reached $\R^{2n}$, which is optimal: the
projective plane $\R P^2$ does not embed in $\R^3$ (a fact we quote and do not prove), although it embeds in $\R^4$ (Exercise~\ref{ex:rp2}).
\end{remark}

Whitney's theorem closes the circle that Riemann opened. Riemann wanted geometry intrinsic, read from inside, without an ambient space; Weyl and Whitney
made that precise with charts; and Whitney then showed that every intrinsic manifold could, after all, be placed in an ambient space. What is gained by
the abstract point of view is not more manifolds but more freedom: we can now work with $\R P^n$, with groups of matrices and with the configuration
spaces of robot arms without first choosing how to draw them.

\begin{figure}[ht]
\centering
\fig{ch01-arm}
\caption{A two-link arm with $\ell_1=1$ and $\ell_2=0.6$, anchored at the origin, reaches the point $w$ in two ways, elbow up and elbow down. Its
workspace is the shaded annulus $\ell_1-\ell_2\le|w|\le\ell_1+\ell_2$. The two boundary circles are exactly the critical values of the arm map of
Exercise~\ref{ex:twolinkarm}: there, the two solutions merge and the arm loses a direction of motion.}
\label{fig:arm}
\end{figure}

\section{Exercises}

\begin{exercise}[$\star$]\label{ex:cubeinverse}
\emph{The hypotheses of the inverse function theorem.} (a) Show that $f(x)=x^3$ is a smooth bijection of $\R$ whose inverse is not differentiable at $0$.
(b) Let $f(x,y)=(e^x\cos y,e^x\sin y)$. Show that $Df(x,y)$ is invertible everywhere but $f$ is not injective. What does the inverse function theorem
give for $f$, and how is $f$ related to the complex exponential?
\end{exercise}

\begin{exercise}[$\star$]\label{ex:matrixgroups}
\emph{More matrix groups.} (a) Show that $\mathrm{SL}_n(\R)=\{\det=1\}$ is a submanifold of $M_n(\R)$ of dimension $n^2-1$. (Differentiate
$t\mapsto\det((1+t)X)$.) (b) Show that $\mathrm U(n)=\{X\in M_n(\C):X^*X=I\}$ is a submanifold of $M_n(\C)=\R^{2n^2}$ of dimension $n^2$. (c) Show that
$\det\colon\mathrm U(n)\to S^1$ is a submersion and deduce that $\mathrm{SU}(n)$ is a submanifold of dimension $n^2-1$. (d) Identify the tangent spaces at
$I$ of $\mathrm{SL}_n(\R)$, $\mathrm U(n)$ and $\mathrm{SU}(n)$.
\end{exercise}

\begin{exercise}[$\star$]\label{ex:twolinkarm}
\emph{The two-link arm.} Let $\ell_1>\ell_2>0$ and $E\colon T^2\to\R^2=\C$, $E(\theta_1,\theta_2)=\ell_1e^{i\theta_1}+\ell_2e^{i(\theta_1+\theta_2)}$, the
position of the hand of an arm whose shoulder turns by $\theta_1$ and elbow by $\theta_2$. (a) Show that the critical points of $E$ are the
configurations with $\theta_2\in\{0,\pi\}$. (b) Show that the critical values are the two circles $|w|=\ell_1\pm\ell_2$. (c) Show that every $w$ with
$\ell_1-\ell_2<|w|<\ell_1+\ell_2$ has exactly two preimages, every $w$ on the circles exactly one, and every other $w$ none (Figure~\ref{fig:arm}).
\end{exercise}

\begin{exercise}[$\star$]\label{ex:nosubmanifold}
\emph{Two non-examples.} (a) Show that the cone $\{(x,y,z):x^2+y^2=z^2\}$ is not a submanifold of $\R^3$ of any dimension. (Look at curves through $0$.)
(b) Let $\psi(t)=(\sin2t,\sin t)$. Show that $\psi$ is an injective immersion of $\left]0,2\pi\right[$ into $\R^2$, but that its image is not a submanifold.
Why does this not contradict Theorem~\ref{thm:fourdesc}?
\end{exercise}

\begin{exercise}[$\star$]\label{ex:stereographic}
\emph{Stereographic projection.} Let $N=(0,\dots,0,1)$ and $S=-N$ in $S^n\subset\R^{n+1}$, and write $x=(x',x_{n+1})$. Let
$\sigma_N(x)=\frac{x'}{1-x_{n+1}}$ on $S^n\setminus\{N\}$ and $\sigma_S(x)=\frac{x'}{1+x_{n+1}}$ on $S^n\setminus\{S\}$. (a) Show that
$\sigma_N^{-1}(y)=\frac{(2y,|y|^2-1)}{|y|^2+1}$. (b) Show that $\sigma_S\circ\sigma_N^{-1}(y)=\frac y{|y|^2}$ on $\R^n\setminus\{0\}$, so the two charts form an
atlas. (c) Show that this atlas defines the same smooth structure as Proposition~\ref{prop:subisman}.
\end{exercise}

\begin{exercise}[$\star$]\label{ex:minisard}
\emph{Curves do not fill the plane.} Let $U\subset\R^m$ be open and $f\colon U\to\R^n$ a $C^1$ map, with $m<n$. Show that $f(U)$ has measure zero. Compare
with Lebesgue's continuous surjection of $[0,1]$ onto the square (Volume~6, Exercise~4.7).
\end{exercise}

\begin{exercise}[$\star\star$]\label{ex:rp2}
\emph{The projective plane in $\R^4$.} For $v=(x,y,z)\in S^2$ let $\tilde F(v)=(x^2-y^2,\,xy,\,xz,\,yz)$. (a) Show that $\tilde F$ induces a smooth map
$F\colon\R P^2\to\R^4$. (b) Show that $F$ is injective. (Compute $x^2+y^2$ from $F$.) (c) Show that $F$ is an immersion, and conclude that $\R P^2$ is
diffeomorphic to a compact surface in $\R^4$.
\end{exercise}

\section{Hints}

\begin{hint}{ex:cubeinverse}
(a) Compute the difference quotient of $x^{1/3}$ at $0$. (b) $\det Df=e^{2x}$, and $f(x,y+2\pi)=f(x,y)$.
\end{hint}
\begin{hint}{ex:matrixgroups}
(b) Copy Example~\ref{ex:on} with Hermitian matrices, of real dimension $n^2$. (c) $t\mapsto e^{it}X$ is a curve in $\mathrm U(n)$.
\end{hint}
\begin{hint}{ex:twolinkarm}
Two vectors $a,b\in\C$ are dependent over $\R$ if and only if $\operatorname{Im}(a\bar b)=0$. For (c), $|w|^2=\ell_1^2+\ell_2^2+2\ell_1\ell_2\cos\theta_2$.
\end{hint}
\begin{hint}{ex:nosubmanifold}
Use Proposition~\ref{prop:tangentsub}: find curves through $0$ whose velocities span too much. For (b), the formula for $\psi$ also makes sense near $t=0$.
\end{hint}
\begin{hint}{ex:stereographic}
For (c), $\sigma_N$ is the restriction of a smooth map on $\{x_{n+1}\ne1\}$, so the chain rule applied to $\sigma_N\circ\sigma_N^{-1}=\mathrm{id}$ shows
$\sigma_N^{-1}$ is an immersion.
\end{hint}
\begin{hint}{ex:minisard}
Extend $f$ to $U\times\R^{n-m}$ without using the new variables.
\end{hint}
\begin{hint}{ex:rp2}
(b) $(x^2+y^2)^2=(x^2-y^2)^2+4x^2y^2$. (c) Show that if $h\perp v$ and $D\tilde F(v)h=0$, then $hv^\top+vh^\top$ is diagonal; apply it to $v$.
\end{hint}

\section{Solutions}

\begin{solution}{ex:cubeinverse}
\idea Each part removes one hypothesis and watches the conclusion fail: first invertibility of the derivative, then the word ``locally''.
\writeup (a) $f$ is smooth and strictly increasing with limits $\pm\infty$, so it is a bijection, with inverse $g(y)=y^{1/3}$ (the real cube root). At $0$,
$\frac{g(y)-g(0)}y=y^{-2/3}\to+\infty$, so $g$ is not differentiable at $0$. The inverse function theorem does not apply there, since $f'(0)=0$. If $g$ were
differentiable at $0$, the chain rule applied to $f\circ g=\mathrm{id}$ would give $f'(0)g'(0)=1$, impossible: a smooth map with a singular derivative
never has a differentiable inverse. (b) $Df=\bigl(\begin{smallmatrix}e^x\cos y&-e^x\sin y\\ e^x\sin y&e^x\cos y\end{smallmatrix}\bigr)$, with determinant
$e^{2x}>0$. But $f(x,y+2\pi)=f(x,y)$. In complex notation $f(x,y)=e^{x+iy}$, the complex exponential. The theorem gives local inverses: these are the
branches of the logarithm, each defined on a small disc avoiding $0$. A global inverse does not exist, for the same reason that the logarithm has no
continuous branch on $\C\setminus\{0\}$ (Volume~7, Chapter~9). Being locally invertible everywhere is a statement about each point, and injectivity is a
statement about pairs of points; no amount of local information gives the second.
\end{solution}

\begin{solution}{ex:matrixgroups}
\idea A group defined by an equation $\Phi(X)=c$ is a manifold as soon as the derivative of $\Phi$ is onto at every solution, and the group structure
helps produce a preimage: one only needs to move along a curve that changes $\Phi$ in a known way.
\writeup (a) $\det((1+t)X)=(1+t)^n\det X$, so for $\det X=1$ the derivative of $\det$ at $X$ in the direction $X$ is $n\ne0$. A non-zero linear form on
$M_n(\R)$ is surjective onto $\R$, so by (b) of Theorem~\ref{thm:fourdesc} $\mathrm{SL}_n(\R)$ is a submanifold of dimension $n^2-1$.

(b) Let $\mathrm{Herm}_n$ be the real vector space of Hermitian matrices; a Hermitian matrix is fixed by $n$ real diagonal entries and $\frac{n(n-1)}2$
complex entries above the diagonal, so $\dim_\R\mathrm{Herm}_n=n^2$. Let $F(X)=X^*X-I\in\mathrm{Herm}_n$. Then $DF(X)H=H^*X+X^*H$, and for $X\in\mathrm U(n)$,
$S\in\mathrm{Herm}_n$, the matrix $H=\frac12XS$ gives $H^*X+X^*H=\frac12SX^*X+\frac12X^*XS=S$. So $\mathrm U(n)$ is a submanifold of dimension $2n^2-n^2=n^2$.

(c) The determinant of a unitary matrix has modulus $1$, so $\det$ maps $\mathrm U(n)$ smoothly to $S^1$ (Proposition~\ref{prop:subisman}). For
$X\in\mathrm U(n)$ the curve $t\mapsto e^{it}X$ lies in $\mathrm U(n)$ and $\det(e^{it}X)=e^{int}\det X$ has derivative $in\det X\ne0$ at $t=0$. So the image
of $d\det_X$ contains a non-zero vector of the line $T_{\det X}S^1$, and $\det$ is a submersion. By Theorem~\ref{thm:regularvalue},
$\mathrm{SU}(n)=\det^{-1}(1)$ is an embedded submanifold of $\mathrm U(n)$ of dimension $n^2-1$; composing the charts, it is also a submanifold of
$M_n(\C)$.

(d) By Proposition~\ref{prop:tangentsub} and Theorem~\ref{thm:regularvalue}, tangent spaces are kernels. The derivative of $\det$ at $I$ is the trace
(expand $\det(I+tH)=1+t\operatorname{tr}H+O(t^2)$), so $T_I\mathrm{SL}_n(\R)=\{\operatorname{tr}H=0\}$. $T_I\mathrm U(n)=\{H^*+H=0\}$, the skew-Hermitian
matrices, and $T_I\mathrm{SU}(n)$ is the set of skew-Hermitian matrices of trace $0$. These are the Lie algebras $\mathfrak{sl}_n$, $\mathfrak u(n)$,
$\mathfrak{su}(n)$ of Chapter~7.
\end{solution}

\begin{solution}{ex:twolinkarm}
\idea The derivative of $E$ has two columns, the velocities of the hand when one joint moves; they are dependent exactly when the arm is straight or
folded, because then both joints move the hand perpendicular to the same line.
\writeup (a) In the coordinates $(\theta_1,\theta_2)$ (local charts of $T^2$), $\partial_{\theta_1}E=iE$ and $\partial_{\theta_2}E=i\ell_2e^{i(\theta_1+\theta_2)}$.
With $a=iE$ and $b=i\ell_2e^{i(\theta_1+\theta_2)}$,
\[
a\bar b=E\,\ell_2e^{-i(\theta_1+\theta_2)}=\ell_2(\ell_1e^{-i\theta_2}+\ell_2),
\]
whose imaginary part is $-\ell_1\ell_2\sin\theta_2$. The real $2\times2$ Jacobian is singular if and only if this vanishes, that is, $\theta_2\in\{0,\pi\}$.

(b) $|E|=|\ell_1+\ell_2e^{i\theta_2}|$, which is $\ell_1+\ell_2$ when $\theta_2=0$ and $\ell_1-\ell_2$ when $\theta_2=\pi$. As $\theta_1$ varies the critical
values sweep these two circles.

(c) Since $|E|^2=\ell_1^2+\ell_2^2+2\ell_1\ell_2\cos\theta_2$, a preimage of $w$ must have $\cos\theta_2=c=\frac{|w|^2-\ell_1^2-\ell_2^2}{2\ell_1\ell_2}$. This has
two solutions $\pm\theta_2$ if $|c|<1$, which is equivalent to $\ell_1-\ell_2<|w|<\ell_1+\ell_2$; one if $c=\pm1$; none otherwise. For each such $\theta_2$
the complex number $u=\ell_1+\ell_2e^{i\theta_2}$ is non-zero (as $\ell_1>\ell_2$) and $|u|=|w|$, so $e^{i\theta_1}=\frac wu$ determines $\theta_1$ uniquely
modulo $2\pi$. Hence two, one or no preimages. The two preimages are the elbow-up and elbow-down postures; as $w$ approaches a boundary circle they merge,
and at the merger the arm can no longer move its hand radially: the critical points are the singular configurations of robotics. Their images are two
circles, of measure zero, as Corollary~\ref{cor:sardmfd} predicts.
\end{solution}

\begin{solution}{ex:nosubmanifold}
\idea If a set is a $k$-dimensional submanifold, its tangent space at a point is $k$-dimensional, and every curve in the set has its velocity there.
So it is enough to exhibit curves whose velocities span more than the set can afford.
\writeup (a) Suppose the cone $K$ is a $k$-dimensional submanifold near $0$. For each angle $\alpha$, the curve $t\mapsto(t\cos\alpha,t\sin\alpha,t)$ lies on $K$, so
$(\cos\alpha,\sin\alpha,1)\in T_0K$. Three such vectors (for $\alpha=0,\frac\pi2,\pi$) are independent, so $k=3$, and a $3$-dimensional submanifold of $\R^3$ is
an open set (in (a) of Theorem~\ref{thm:fourdesc}, $\varphi(K\cap U)=V$, so $K\cap U=U$). But no ball around $0$ lies in $K$: it contains $(0,0,t)$ for small
$t\ne0$, and $0\ne t^2$.

(b) $\psi'(t)=(2\cos2t,\cos t)$ vanishes only if $\cos t=0$ and $\cos2t=2\cos^2t-1=0$, impossible. If $\psi(t)=\psi(s)$, then $\sin s=\sin t$, so $s=t$ or
$s=\pi-t$ modulo $2\pi$; in the second case $\sin2s=-\sin2t$, which forces $\sin2t=0$, and the only candidate $t=\pi$ would need $s\in\{0,2\pi\}$, outside the
interval. So $\psi$ is an injective immersion. Let $S=\psi(\left]0,2\pi\right[)$. The curve $c(t)=(\sin2t,\sin t)$ for $|t|<\varepsilon$ lies in $S$ (for
$t<0$, $c(t)=\psi(t+2\pi)$; for $t>0$, $c(t)=\psi(t)$; and $c(0)=\psi(\pi)$), with velocity $(2,1)$, and $t\mapsto\psi(\pi+t)$ has velocity $(2,-1)$.
As in (a), if $S$ were a submanifold it would be $2$-dimensional, hence open in $\R^2$ near $0$, but $S$ is the image of a $C^1$ curve, which has measure
zero (Exercise~\ref{ex:minisard}). No contradiction with Theorem~\ref{thm:fourdesc}: $\psi$ is not a homeomorphism onto $S$, since
$\psi(\frac1k)\to\psi(\pi)$ while $\frac1k\not\to\pi$. The branch arriving at the origin along $t\to0$ is ``far'' in the parameter and ``near'' in the image.
\end{solution}

\begin{solution}{ex:stereographic}
\idea Stereographic projection is the line through the pole and the point, cut by the equatorial plane; inverting it is solving a quadratic that has one
known root, the pole.
\writeup (a) Let $x=\frac{(2y,|y|^2-1)}{|y|^2+1}$. Then $|x|^2=\frac{4|y|^2+(|y|^2-1)^2}{(|y|^2+1)^2}=1$, $x\ne N$, and $1-x_{n+1}=\frac2{|y|^2+1}$, so
$\sigma_N(x)=y$. Conversely $\sigma_N$ is injective on $S^n\setminus\{N\}$: from $y=\sigma_N(x)$ one computes $|y|^2=\frac{1-x_{n+1}^2}{(1-x_{n+1})^2}=\frac{1+x_{n+1}}{1-x_{n+1}}$, which
determines $x_{n+1}$, and then $x'=(1-x_{n+1})y$. So $\sigma_N$ is a bijection onto $\R^n$ with the stated inverse, both continuous: a chart. (b) For
$y\ne0$, $1+x_{n+1}=\frac{2|y|^2}{|y|^2+1}$ and $x'=\frac{2y}{|y|^2+1}$, so $\sigma_S(x)=\frac y{|y|^2}$, a smooth map of $\R^n\setminus\{0\}$ that is its own inverse. By symmetry
the same holds the other way. (c) $\sigma_N$ is the restriction to $S^n$ of the smooth map $x\mapsto\frac{x'}{1-x_{n+1}}$ on $\{x_{n+1}\ne1\}$, so it is smooth
for the submanifold structure (Proposition~\ref{prop:subisman}). Its inverse $\sigma_N^{-1}$ is smooth from $\R^n$ to $\R^{n+1}$, and from
$\sigma_N\circ\sigma_N^{-1}=\mathrm{id}$ the chain rule gives that $D\sigma_N^{-1}(y)$ is injective; together with (a), $\sigma_N^{-1}$ is a parametrization in the
sense of (d). In a straightening chart $\varphi$, the map $\pi\circ\varphi\circ\sigma_N^{-1}$ is smooth, and its inverse $\sigma_N\circ\varphi^{-1}(\cdot,0)$ is
smooth too; so the stereographic charts are compatible with the charts of Proposition~\ref{prop:subisman}, and the two atlases lie in the same maximal atlas.
\end{solution}

\begin{solution}{ex:minisard}
\idea A map from a lower dimension is a map from the full dimension that ignores some variables; then every point is critical, and the null-set lemma
does the rest.
\writeup Let $G(x,z)=f(x)$ on the open set $U\times\R^{n-m}\subset\R^n$. It is $C^1$, and $f(U)=G(U\times\{0\})$. The set $U\times\{0\}$ has measure zero in
$\R^n$ (it lies in a hyperplane), so $f(U)$ is null by Lemma~\ref{lem:nullsets}. Lebesgue's curve is continuous but not $C^1$, nor even Lipschitz: the
proof of the lemma only needs a bound $|f(x)-f(y)|\le L|x-y|$, which turns a cover by small cubes into a cover by cubes of comparable size, and the same
proof shows that a Lipschitz curve has a null image. A space-filling curve must therefore stretch small intervals into comparatively huge sets, and
this is exactly what continuity alone allows.
\end{solution}

\begin{solution}{ex:rp2}
\idea The map remembers the products of pairs of coordinates, which is almost the matrix $vv^\top$, and that matrix determines the line. For the immersion
one uses the same idea infinitesimally: the derivative of $v\mapsto vv^\top$ is $h\mapsto hv^\top+vh^\top$, which cannot vanish on $h\perp v$.
\writeup (a) $\tilde F(-v)=\tilde F(v)$, since each component is quadratic. The quotient map $\pi\colon S^2\to\R P^2$ restricts, on each open hemisphere
where one coordinate is positive, to a diffeomorphism onto a chart domain of Example~\ref{ex:rpn}: on $\{z>0\}$, followed by the chart
$[x:y:z]\mapsto(\frac xz,\frac yz)$, it reads $(x,y,z)\mapsto(\frac xz,\frac yz)$, with smooth inverse $(u,s)\mapsto\frac{(u,s,1)}{\sqrt{1+u^2+s^2}}$. So
$F=\tilde F\circ(\pi|_{\text{hemisphere}})^{-1}$ is well defined and smooth on each chart domain.

(b) Let $(a,b,c,d)=\tilde F(v)$ with $|v|=1$. Since $(x^2+y^2)^2=(x^2-y^2)^2+4x^2y^2=a^2+4b^2$, we get $x^2+y^2=\sqrt{a^2+4b^2}$, hence
$x^2=\frac{\sqrt{a^2+4b^2}+a}2$, $y^2=\frac{\sqrt{a^2+4b^2}-a}2$ and $z^2=1-x^2-y^2$. Together with $xy=b$, $xz=c$, $yz=d$, all the entries of $vv^\top$ are
determined by $F([v])$. The matrix $vv^\top$ is the projection onto the line $\R v$, so $[v]$ is determined.

(c) Since $\pi$ is a local diffeomorphism, $F$ is an immersion if $D\tilde F(v)$ is injective on $T_vS^2=v^\perp$. Let $h\perp v$ with $D\tilde F(v)h=0$:
\[
2xh_1-2yh_2=0,\quad yh_1+xh_2=0,\quad zh_1+xh_3=0,\quad zh_2+yh_3=0.
\]
The last three say that the symmetric matrix $P=hv^\top+vh^\top$ has zero off-diagonal entries; its diagonal is $(2xh_1,2yh_2,2zh_3)$. Put $s=xh_1=yh_2$.
From $\langle v,h\rangle=0$, $zh_3=-2s$, so $P=\operatorname{diag}(2s,2s,-4s)$. On the other hand $Pv=h|v|^2+v\langle h,v\rangle=h$. Hence
$h=(2sx,2sy,-4sz)$, and then $s=xh_1=2sx^2$ and $s=yh_2=2sy^2$. If $s\ne0$ this gives $x^2=y^2=\frac12$, so $z=0$ and $-2s=zh_3=0$, a contradiction. So
$s=0$, $P=0$ and $h=Pv=0$. Finally $\R P^2=\pi(S^2)$ is compact, so by Proposition~\ref{prop:embedding} $F$ is an embedding and $\R P^2$ is diffeomorphic
to the compact surface $F(\R P^2)\subset\R^4$.
\end{solution}

\section*{Sources and further reading}
\addcontentsline{toc}{section}{Sources and further reading}

B.~Riemann, ``\"Uber die Hypothesen, welche der Geometrie zu Grunde liegen'', \emph{Abh.\ K\"onigl.\ Ges.\ Wiss.\ G\"ottingen} 13 (1868), 133--152,
published by Dedekind from Riemann's manuscript; W.~K.~Clifford's English translation appeared in \emph{Nature} 8 (1873). The account of the choice of
topic and of Gauss's reaction is Dedekind's, in the biography that closes Riemann's \emph{Gesammelte mathematische Werke} (Teubner, 1876); we have
paraphrased it. D.~Laugwitz, \emph{Bernhard Riemann 1826--1866: Turning Points in the Conception of Mathematics} (Birkh\"auser, 1999), tells the
story in full and discusses the lecture.

H.~Poincar\'e, ``Analysis situs'', \emph{J.\ \'Ec.\ Polytech.} (2) 1 (1895), 1--121. H.~Weyl, \emph{Die Idee der Riemannschen Fl\"ache} (Teubner, 1913).
H.~Whitney, ``Differentiable manifolds'', \emph{Ann.\ of Math.} 37 (1936), 645--680, and ``The self-intersections of a smooth $n$-manifold in $2n$-space'',
\emph{Ann.\ of Math.} 45 (1944), 220--246. A.~Sard, ``The measure of the critical values of differentiable maps'', \emph{Bull.\ Amer.\ Math.\ Soc.} 48
(1942), 883--890; A.~P.~Morse, ``The behavior of a function on its critical set'', \emph{Ann.\ of Math.} 40 (1939), 62--70; H.~Whitney, ``A function not
constant on a connected set of critical points'', \emph{Duke Math.\ J.} 1 (1935), 514--517.

Our treatment of the inverse function theorem follows W.~Rudin, \emph{Principles of Mathematical Analysis}, 3rd ed. (McGraw-Hill, 1976), Chapter~9; the
four descriptions of submanifolds and the tangent space follow J.~Lafontaine, \emph{An Introduction to Differential Manifolds} (Springer, 2015), and the
abstract theory J.~M.~Lee, \emph{Introduction to Smooth Manifolds}, 2nd ed. (Springer, 2013), Chapters~1--5, where the general Sard theorem and the
strong Whitney embedding theorem are also proved. J.~Milnor, \emph{Topology from the Differentiable Viewpoint} (University Press of Virginia, 1965), proves
Sard's theorem in forty lines and is the model for everything in Chapter~3. For the robotics, J.~J.~Craig, \emph{Introduction to Robotics: Mechanics and
Control}, 3rd ed. (Pearson, 2005), Chapters~4--5, discusses workspaces and singular configurations of arms.
